\chapter{The Tug-of-War Behind Periodicity}

\begin{kaichang}
Stir a small spoonful of salt into clear water. It disappears, but a taste confirms it remains. Earlier chapters explained that salt is an orderly crystal of sodium ions \chem{Na^+} and chloride ions \chem{Cl^-}. In water they separate and swim independently.

Now make two guesses. First, when sodium \chem{Na} loses an electron to become an ion, does it grow or shrink? Second, when chlorine \chem{Cl} gains an electron, does it become fatter or thinner?

Many people instinctively say losing something makes you smaller and gaining something larger. That is half right, but the surprising part is the scale: sodium really shrinks, far more than you might expect, with radius nearly halved; chloride grows comparably dramatically. Stranger still, the same mechanics explains both changes.

This mechanics has no formal name, so let us call it a \textbf{tug-of-war}: the nucleus pulls electrons inward while electrons push one another outward. Chapter 10 outlined the periodic-table map. Here we derive its leftward and rightward trend arrows one by one from that contest. You will find that periodicity needs no memorization; it records two competing forces.
\end{kaichang}

\section{Measuring an Atom's Waist}

Before discussing the contest, we need its arena's size. How large is an atom?

That is much harder than measuring your waist. A tape can follow skin's clear boundary. Atoms have none. Chapter 9 explained that electrons are probability clouds, not planets following fixed orbits. The cloud thins outward but never reaches zero. There is no shell or line marking an end.

Where do textbook radii come from, then? A practical solution measures \textbf{distances between atoms}, not isolated atoms. In metallic sodium, X-ray diffraction infers neighboring nuclear separations from wave-interference patterns in a regular array; half the distance becomes an atomic radius. In \chem{Cl_2}, half the chlorine--chlorine nuclear distance gives another radius. Different assumptions produce somewhat different numbers. Strictly, radius is not absolute: it measures how far an atom keeps from its neighbors under a particular relationship.

Translate distance into intuition. An atomic radius is about 100 picometers, pm, or $10^{-10}$ meters. A hundred million atoms side by side extend roughly 1 centimeter. Your kitchen foil is about 0.016 millimeters thick, seemingly negligible, but dividing by aluminum's approximate diameter, $2.9\times 10^{-10}$ meters, gives \textbf{over fifty thousand atomic layers}.

Now examine measured radii for some second- and third-period elements, in pm. Methods differ slightly; these are common values:

\begin{table}[htbp]
\centering
\begin{tabular}{lccccccc}
\toprule
Period 2 & Li & Be & B & C & N & O & F \\
\midrule
Radius/pm & 128 & 96 & 84 & 76 & 71 & 66 & 57 \\
\bottomrule
\end{tabular}
\end{table}

\begin{table}[htbp]
\centering
\begin{tabular}{lccccccc}
\toprule
Period 3 & Na & Mg & Al & Si & P & S & Cl \\
\midrule
Radius/pm & 166 & 141 & 121 & 111 & 107 & 105 & 102 \\
\bottomrule
\end{tabular}
\end{table}

The pattern is immediate: \textbf{atoms shrink from left to right across a row}. Compare lithium at 128, sodium at 166, and potassium a little above 200 pm vertically: \textbf{atoms grow down a column}.

Odd, is it not? Moving right adds electrons, yet atoms shrink rather than expand. Increasing positive nuclear charge seems more influential than increasing electron count. That is our first clue. Let us meet the contestants.

\section{Two Forces, One Arena}

Zoom into an atom and inventory its forces:

\begin{itemize}[itemsep=2pt]
\item \textbf{Contestant one: nuclear attraction.} Positive nuclei attract negative electrons. This force has a crucial trait: \textbf{extreme sensitivity to distance}. Double separation and attraction falls to one quarter, an inverse-square law; halve it and attraction quadruples. Low-floor electrons are held tightly, while those higher up strain the short-armed nucleus's reach.
\item \textbf{Contestant two: electron repulsion.} All electrons are negative, so like charges repel. Each feels an inward nuclear pull and outward pushes from surrounding electrons. More electrons mean more internal squabbling.
\end{itemize}

An atom's build and personality depend on the score between these contestants. Its most important concept is the \textbf{shielding effect}.

Imagine an outer electron looking toward the nucleus through inner electrons. Those negative charges form a net between it and the positive nucleus, \textbf{blocking part of the attraction}. The outer electron feels a discounted net pull instead of the full nuclear charge. That discounted number is \textbf{effective nuclear charge}, commonly $Z^{*}$.

A simple analogy can help you feel these forces:

\begin{shiyan}
\textbf{Electrostatic Attraction and Shielding}

\textbf{Materials}: a plastic ruler or balloon, dry hair or a sweater, paper scraps torn to rice-grain size, and kitchen aluminum foil.

\textbf{Procedure}:
\begin{enumerate}[itemsep=1pt]
\item Scatter the scraps on a dry table.
\item Rub the ruler rapidly against hair or sweater a dozen or more times, then approach the scraps slowly without touching. Record the distance at which they begin jumping up.
\item Repeat at different distances to appreciate that attraction works best close up.
\item Have a partner hold foil upright between ruler and scraps without touching the ruler. Bring the charged ruler close again. Do the scraps move? Remove the foil and repeat.
\end{enumerate}

\textbf{What to observe}: A closer ruler picks scraps up more readily; slightly farther away, nothing moves. Attraction weakens rapidly with distance. Thin foil between them makes the scraps suddenly play dead, as though attraction switched off.

\textbf{Safety}: Generate charge only by rubbing hair or sweaters. \textbf{Never} play with static near outlets or wires, and never insert metal objects into sockets. Foil edges are sharp; avoid cuts.

\textbf{What this models}: The ruler's charge resembles the nucleus's positive charge: its pull weakens with distance, contestant one's personality. Foil blocking the field is a macroscopic version of shielding. Be clear, however: atomic shielding involves quantum repulsion and electron-cloud arrangements, only \emph{analogous} to foil screening an electric field, not the same mechanism. The analogy builds intuition and nothing more.
\end{shiyan}

\section{A Mental Scorecard: Effective Nuclear Charge}

Turn shielding into a rough mental model with only two rules:

\begin{qiaoliang}
\textbf{Mental Rules for Effective Nuclear Charge: Rough but Useful}

Divide electrons into inner and outer shells. Assume each inner electron cancels roughly one nuclear charge, while same-shell electrons crowd together and barely block one another's view of the nucleus.

Thus $Z^{*} \approx$ nuclear charge $-$ inner-electron count, approximately the outer-electron count.

Some examples:
\begin{itemize}[itemsep=1pt]
\item Lithium (\chem{Li}, nuclear charge +3, configuration 2, 1): two inner electrons screen two charges, leaving its outer electron a net pull around $3-2=1$.
\item Sodium (\chem{Na}, +11, configuration 2, 8, 1): $11-10=1$. Potassium (\chem{K}, +19, configuration 2, 8, 8, 1): $19-18=1$. \textbf{Effective nuclear charge is nearly equal within a group.}
\item Fluorine (\chem{F}, +9, configuration 2, 7): each outer electron feels about $9-2=7$. Chlorine (\chem{Cl}, +17, configuration 2, 8, 7): $17-10=7$.
\item Across period 3: sodium about 1, magnesium about 2, aluminum about 3, onward to chlorine about 7. \textbf{Net pull increases step by step across a period.}
\end{itemize}
\end{qiaoliang}

With this scorecard, the two major radius trends immediately acquire causes.

\textbf{Why do atoms shrink across a period?} Electrons enter the \emph{same shell}, so the building gains no floors, while nuclear charge rises step by step. Same-shell electrons provide little shielding, so each electron's net pull $Z^{*}$ rises from 1 to 7. Greater pull tightens the entire cloud. From lithium to fluorine, the atom does not simply deflate; the nucleus wins the tug-of-war.

\textbf{Why do atoms grow down a group?} $Z^{*}$ stays similar, all 1 or all 7, but electrons move up a floor. Attraction is intensely distance-sensitive: higher shells mean greater distance and weaker reach. Distance wins this round.

This is how the score enters the periodic table. Remembering the reasoning matters more than any arrow: \textbf{trends are derived, not memorized}.

\section{The Price of Taking an Electron}

Radius measures arena size. A more direct question is \textbf{how much ransom must you pay to pull an electron out of an atom?}

The ransom is \textbf{ionization energy}, the energy needed to remove the outermost electron from a gaseous atom and make a positive ion. Symbolically:

\[\chem{Na(g)} \rightarrow \chem{Na^{+}(g)} + e^{-} \qquad \Delta E = +496\ \mathrm{kJ/mol}\]

Notice the plus sign: 496 is energy \textbf{you pay}, not a gift from the atom. This echoes the book's recurring account rule: separating mutually attracting objects costs energy; bringing them together releases it. Chapter 17's chemical bonds will use this repeatedly.

Here are first ionization energies for period 2, in kJ/mol:

\begin{table}[htbp]
\centering
\begin{tabular}{lcccccccc}
\toprule
Element & Li & Be & B & C & N & O & F & Ne \\
\midrule
First ionization energy & 520 & 900 & 801 & 1086 & 1402 & 1314 & 1681 & 2081 \\
\bottomrule
\end{tabular}
\end{table}

The overall pattern matches our prediction: increasing $Z^{*}$ tightens electrons and raises ransom across the row, peaking at neon's full shell. Down a group also fits: sodium requires 496, potassium only 419. Potassium's outer electron lives on floor four, farther out and easier to take. This quantifies Chapter 10's higher-floor/easier-loss rule and explains Group 1's increasing temper downward (Chapter 12 examines water reactions).

Did you notice two slips? Beryllium's 900 is \emph{higher} than boron's 801, and nitrogen's 1402 \emph{higher} than oxygen's 1314. Moving right increases $Z^{*}$, yet ransom falls. This is not bad measurement but a finer rule: \textbf{room type and shared occupancy matter too}. The mentally demanding details await the closing challenge. For now, remember that periodicity is a real overall pattern, but \textbf{a trend with exceptions}, themselves explained by deeper energy accounting. That is science at its most fascinating.

Get a feel for 496 kJ/mol. A gram of sodium is about $1/23$ mole, so removing every atom's outer electron requires

\[496\ \mathrm{kJ/mol} \times \frac{1}{23}\ \mathrm{mol} \approx 21.6\ \mathrm{kJ}\]

What is 21.6 kJ? Water's heat capacity is about 4.2 kJ per kilogram per degree Celsius, so this can heat 1 kilogram of water by about $5\,^{\circ}\mathrm{C}$, or power a 60-watt bulb for 6 minutes. One small gram of metal holds that much electron ransom. Nuclear attraction is no trivial matter.

Taking requires a recipient. Some atoms accept an electron and pay energy back rather than charge a fee: a gaseous chlorine atom gaining an electron to become chloride \textbf{releases} 349 kJ/mol. This is \textbf{electron affinity}. Release means the combination of chloride ion and free electron has lower total energy than chlorine atom plus electron. Again, change heads toward lower energy, not an atom's desire for electrons.

Combine the accounts: sodium costs 496 and chlorine returns only 349, leaving a net 147 kJ/mol expense. Why can salt exist stably if the total loses energy credit? Keep that question. An exercise will ask you to calculate it, and a new force in Part III supplies the missing entry.

\section{Electronegativity: The Tug-of-War Scoreboard}

Ionization energy and electron affinity concern complete removal or acceptance. More often, neither atom fully lets go: they \textbf{share} an electron pair to make a chemical bond, Part III's theme. Is sharing fair? Rarely. The shared electrons resemble a ribbon tied at the middle of a tug-of-war rope, almost always pulled toward the stronger side.

Physicists named this ability to pull shared electrons \textbf{electronegativity}. On a human-defined scoreboard, fluorine leads at about 4.0; oxygen is 3.4, chlorine 3.2, nitrogen 3.0, carbon 2.6, hydrogen 2.2, sodium only 0.9, and bottom-ranked cesium about 0.8. You can predict the trend: increasing toward fluorine at upper right, decreasing toward cesium at lower left. The same causes apply: larger $Z^{*}$ and lower floors give the nucleus a stronger hand.

How was this scoreboard established? Nobody can hook electrons to a spring balance. A beautiful evidence story lies behind it.

\begin{zhengju}
In 1932, the American chemist Linus Pauling studied a seemingly minor accounting puzzle: why was the \chem{H-Cl} bond energy noticeably greater than the \textbf{average} of \chem{H-H} and \chem{Cl-Cl} bond energies?

He reasoned that equal hydrogen and chlorine pulls would distribute shared electrons evenly, giving no reason for the mixed bond to exceed the average pure-bond strength. An unbiased rope gives neither side an advantage. Yet experiments repeatedly found unlike-atom bonds \textbf{stronger} than the average. Where did the extra strength come from? Pauling inferred uneven electron sharing. Pulling electrons toward one end made it slightly negative and the other slightly positive; electrostatic attraction added a bonus strengthening the bond. Greater imbalance gave a greater bonus.

He then did something unusual for chemists: converted the excess bond energy into a ranking from 0 to 4, the electronegativity scale of electron-pulling power. The entire table came from calorimetry, without seeing a single electron directly.

The definition was controversial as an \textbf{inferred} indirect quantity rather than something weighed on a balance. Alternatives followed. Mulliken proposed averaging electron-taking and electron-retaining abilities, electron affinity and ionization energy. Allred and Rochow estimated directly from effective nuclear charge divided by radius squared. The rankings differed in detail but almost not at all in order: fluorine first, cesium last. Independent calculations converging on the same order suggested electronegativity captured something real. Many useful scientific concepts are like this: not directly readable, yet unavoidable however calculated.
\end{zhengju}

Electronegativity has more uses than you might imagine. Large differences pull electrons strongly to one side, producing ions as in salt; small differences permit sharing, making covalent bonds in water, sugar, and fats. Why does water dissolve salt but not oil? Molecular positive and negative ends created by electronegativity differences hold the answer, developed in Part III. Mass spectrometers in hospital and anti-doping laboratories likewise first use ionization energy to turn sample molecules into charged ions, then sort by mass, detecting banned substances at parts per billion. Neon signs use high voltage to strip electrons from noble gases, making them conductive and luminous. Noble gases avoid ordinary electronegativity tug-of-war, so the tubes are not corroded.

\section{The Trends in Everyday Life}

These scores are no laboratory ornaments. Many familiar facts follow from them.

Why can helium balloons float safely over children at birthday parties? Helium's outer-electron ransom is among the highest, beyond other atoms' budgets, so it avoids reactions. Hydrogen balloons instead burn fiercely near flames, hence strict helium-only rules at large events. Laboratory cesium is sealed in protective-gas-filled glass tubes, visibly soft and golden. Its bottom-ranked ionization energy makes outer electrons almost free, immediately traded upon air exposure; Chapter 12 shows its fiery meeting with water. An elder's gold ring stays bright and unrusted for decades partly because gold's high ionization energy and substantial electronegativity hold electrons tightly, making oxygen's deal difficult. Iron nails rust within months. Chapter 15 accounts for metal differences and rust's attack and defense. Blood potassium and sodium tests measure \chem{K^+} and \chem{Na^+}, since these atoms long ago surrendered their outer electrons in the body: ransom was too cheap to retain them.

Which electrons are easy to take and which hard to disturb writes the whole periodic table's expression in advance through this contest.

\section{Where This Map Reaches Its Limits}

Every useful model fails somewhere. Before using the tug-of-war model, read its warranty conditions.

\begin{itemize}[itemsep=2pt]
\item \textbf{Trends have exceptions.} Beryllium and nitrogen already showed ionization-energy slips. A filled or half-filled room type brings extra stability and ransom. The overall trend is a slope, exceptions steps; both are real.
\item \textbf{The central district is less obedient.} Transition metals fill inner d levels, with complicated shielding and small differences in neighboring radii and ionization energies. Trends almost flatten. This explains their similar personalities and multiple valences; Chapter 15 visits them.
\item \textbf{These numbers are not simply weighed.} Radii depend on settings, differing in metals, molecules, and crystals; electronegativity depends on its defining scheme. New calculations shift values slightly, but \textbf{trends and rankings remain firm}. Comparing magnitudes and directions is reliable; fighting over a decimal place misses the point.
\item \textbf{Ion size also depends on neighbors.} The number of surrounding oppositely charged ions in a crystal affects size slightly. Ionic radius likewise belongs to a particular relationship.
\end{itemize}

The warranty in one sentence: this model is \textbf{scaffolding for reasoning}, not a memorization list. Before applying a trend, ask which sides $Z^{*}$, shell level, and room type support in this particular contest.

\begin{wuqu}
``Electrons are so light that gaining or losing them changes only the charge symbol, hardly the atom's size.''

The mistake is seductive because its first half is true. An electron weighs about 1/1836 of a proton, so sodium loses almost no weight with one electron. But size depends not on electron \emph{mass} but on \emph{the balance of pull and repulsion}. Sodium's radius is about 186 pm, its ion's only 102 pm, roughly 55\% as large. Volume shrinks cubically to about $(102/186)^3 \approx 1/6$! Conversely, chlorine's roughly 102 pm grows to about 181 pm for chloride. A nearly massless electron can multiply volume or reduce it to a fraction. Show these numbers to anyone arguing that light electrons imply similar ionic and atomic sizes.
\end{wuqu}

\section{Back to the Saltwater}

Now we can answer the opening puzzles.

\textbf{Why does sodium shrink so dramatically?} Its configuration is 2, 8, 1. Losing the sole third-shell electron demolishes that entire floor, leaving 10 electrons in shells one and two. Nuclear charge stays +11, but only 10 electrons remain, increasing the net pull each feels and tightening the cloud. Demolition plus tightening lowers radius from about 186 to 102 pm.

\textbf{Why does chlorine grow?} It starts with 2, 8, 7 electrons and nuclear charge +17. Adding the eighteenth electron adds no shell, but the nucleus must hold one more electron and repulsion increases. Effective pull is shared more thinly, loosening the cloud from about 102 to 181 pm.

Your saltwater therefore contains a contrasting pair: small sodium ions and large chloride ions. They do not swim naked, either. Each wears a coat of water molecules. Why do water molecules line up so willingly? Water itself has positive and negative ends, another clue left by electronegativity, revealed when Part III discusses molecular polarity.

\section{The Contest Is Far from Over}

We have unpacked periodicity into nuclear attraction's distance sensitivity and electron repulsion and shielding. Radius, ionization energy, electron affinity, and electronegativity are four readings of one scoreboard.

Next we visit its extreme divisions. Chapter 12 takes the far left, Groups 1 and 2, whose electron surrender is so easy that water can close the deal on contact, with fiery results. Chapter 13 visits the far right: halogens maximize electron-taking, while noble gases withdraw from competition. Even withdrawal is not absolute: xenon has been forced back into play. These extremes perform more dramatically than the middle.

\begin{tiaozhan}
\textbf{Why Does Ionization Energy Slip?} (Skipping this does not interrupt the main story.)

Beryllium's 900 exceeds boron's 801, and nitrogen's 1402 exceeds oxygen's 1314. Room types and sharing explain why. Within a shell, s orbitals approach the nucleus more closely than p orbitals, quantum penetration. Boron's lone 2p electron is therefore easier to remove than beryllium's close-in 2s pair, lowering ransom. Oxygen similarly has three 2p rooms. Electrons first occupy them singly, creating a half-full set, but the fourth must share. Strong same-room repulsion gives oxygen an uncomfortable electron that leaves more easily than nitrogen's singly housed electrons. Filled and half-filled stability is not a mysterious law: \textbf{minimized repulsion and room-type symmetry} jointly balance the account. Detailed shielding calculations, Slater's rules, assign fractional screening to same-shell electrons too, agreeing with our rough model's conclusions while refining the score.
\end{tiaozhan}

\section*{Try It Yourself}

\begin{sikaoti}
\item Estimate outer-electron effective nuclear charge for magnesium (\chem{Mg}, +12) and sulfur (\chem{S}, +16). Which should give electrons, which accept them, and what ionic charges should result?
\item Draw a fluorine atom, +9 nucleus and 9 electrons, beside fluoride, +9 nucleus and 10 electrons. Mark nuclear pulls and electron repulsions, and identify the smaller pull per electron. Note that circles and arrows are representations; electrons do not actually form circles.
\item Rank \chem{Na}, \chem{Na^{+}}, \chem{Cl}, and \chem{Cl^{-}} by decreasing radius, explaining in three sentences using shells, shielding, and repulsion.
\item Plot energy against reaction progress: sodium ionization uphill by $+496$ kJ/mol, then chlorine electron gain downhill by $-349$ kJ/mol, marking net $+147$ kJ/mol. Why can salt crystals remain stable despite this expense? Guess the missing account entry. (Hint: what happens when opposite ions assemble into a crystal?) Chapter 18 checks the answer.
\item Explain to a classmate why beryllium exceeding boron and nitrogen exceeding oxygen in ionization energy do not invalidate periodicity, using the challenge's room-type and sharing language.
\item How much energy ionizes every atom once in 2 grams of sodium, molar mass about 23 g/mol? By how many Celsius degrees could it heat 1 kilogram of water? (Water's heat capacity is about 4.2 kJ/(kg·$^{\circ}$C).)
\end{sikaoti}
